\section{Problem statement}

Existing studies of Nepalese commercial banks have examined associations between
market prices and accounting indicators such as earnings per share, book value
per share, dividend measures, and profitability ratios. Findings vary across
studies, samples, and periods. For example, \citet{chhetri2023} reports a
positive ROA coefficient, \citet{chaudhari2025} reports a negative one, and
\citet{gyawali2022} reports a positive but statistically insignificant ROA
coefficient. For BVPS, \citet{joshi2023} report a positive association while
\citet{pant2024} report a negative, statistically insignificant coefficient.
These estimates are not directly comparable because they use different samples,
periods, variable definitions, and model specifications.

Most Nepalese price-determinant models, whether pooled or panel fixed-effects,
estimate an average accounting-price slope across their observation periods
\citep{chhetri2023,gurungdahal2025,bhattarai2024mps}. They do not explicitly
test whether the association between accounting fundamentals and market
valuation changed across the early-2023 banking consolidation. Four of the five
sampled commercial banks underwent major consolidations that commenced joint
operation in January--February 2023, altering their corporate scale and
reporting boundaries
\citep{nrb2023list}.

A second design consideration is the use of price-derived ratios such as P/E as
explanatory variables. Because $\text{P/E}=\text{MPS}/\text{EPS}$, including
P/E in a model whose outcome is MPS creates definitional overlap and
complicates interpretation. This proposal excludes P/E and defines the
explanatory variables using accounting information that is not algebraically
constructed from the observed dependent-variable price.

A third design consideration is timing. This proposal defines the dependent
variable as the market price per share in a standardized window on or after the
statutory nine-month audited-reporting deadline in Section 59(2) of the Banks
and Financial Institutions Act, operationalized at the end of Chaitra
\citep{bafia2017}. The rule makes the observation date reproducible; it does
not establish that each bank released its audited statements by that date.
Where verifiable, actual release dates will be recorded, and this limitation
will constrain interpretation of the price as reflecting the audited accounts.

The study addresses the following questions:
{\renewcommand{\theenumi}{\Roman{enumi}}
\begin{enumerate}[label=\Roman*.,align=right,leftmargin=44.2pt,labelwidth=30pt,labelsep=14.2pt]
  \item What conditional associations do EPS and BVPS have with the selected banks' market prices measured in the reporting-deadline window?
  \item Do the estimated EPS--price and BVPS--price associations differ between the pre- and post-consolidation periods?
  \item Are the principal findings robust to alternative accounting-variable specifications and to the transition-year and entity-timing observations?
\end{enumerate}
}

\subsection{Scope of the study}
\begin{sloppypar}
The study is confined to five purposively selected Nepalese commercial bank
series: GBIME, HBL, KBL, NIMB, and NMB, observed over FY 2076/77--2081/82. It
examines the conditional associations of EPS and BVPS with a ShareSansar-derived
reported share-price close measured in the first eligible trading session on
or after Chaitra-end, subject to the five-calendar-day rule described in the
methodology. Its period comparison tests whether the estimated EPS--price and
BVPS--price slopes differ between the defined pre- and post-consolidation
periods. ROA and eligible DPS are supplementary specifications. The purposive
sample, aggregator-based price measure, reporting-date proxy, and changing
corporate boundaries limit generalization and do not support causal claims
about merger effects.
\end{sloppypar}
